\documentclass[11pt,letterpaper]{article}
\usepackage[T1]{fontenc}
\usepackage[utf8]{inputenc}
\usepackage{lmodern,microtype}
\usepackage[margin=1in,top=.93in,bottom=.93in,headheight=15pt]{geometry}
\usepackage{amsmath,amssymb,amsthm,mathtools,booktabs,longtable,array,enumitem}
\usepackage[table]{xcolor}
\usepackage{fancyhdr}
\definecolor{ink}{HTML}{184D58}
\definecolor{faintink}{HTML}{566C70}
\usepackage[unicode,colorlinks=true,linkcolor=ink,citecolor=ink,urlcolor=ink,breaklinks=true]{hyperref}
\usepackage{bookmark}
\setlength{\emergencystretch}{3em}
\setcounter{tocdepth}{1}
\setlist{nosep,topsep=4pt}
\allowdisplaybreaks[2]
\pagestyle{fancy}\fancyhf{}\fancyhead[L]{\small\color{faintink}The M\"obius function in iPRA}
\fancyhead[R]{\small\color{faintink}Constructive number theory}
\fancyfoot[C]{\small\thepage}
\renewcommand{\headrulewidth}{.35pt}
\newtheorem{theorem}{Theorem}[section]
\newtheorem{lemma}[theorem]{Lemma}
\newtheorem{proposition}[theorem]{Proposition}
\newtheorem{corollary}[theorem]{Corollary}
\theoremstyle{definition}\newtheorem{definition}[theorem]{Definition}
\theoremstyle{remark}\newtheorem{remark}[theorem]{Remark}
\newcommand{\ipra}{\mathrm{iPRA}}
\newcommand{\NN}{\mathbb N}
\newcommand{\ZZ}{\mathbb Z}
\hypersetup{pdftitle={The Möbius function in intuitionistic primitive recursive arithmetic}}
\begin{document}
\thispagestyle{plain}
\noindent{\color{ink}\sffamily\small CONSTRUCTIVE NUMBER THEORY}\par\vspace{8pt}
{\noindent\LARGE\bfseries The M\"obius function in intuitionistic\\[3pt] primitive recursive arithmetic\par}
\vspace{8pt}
\noindent{\small November 22, 2026}\par
\vspace{9pt}
\noindent
The finite sum $\sum_{d\mid n}\mu(d)=[n=1]$ accounts for much of the arithmetic of $\mu$.
It gives inversion, counts of irreducible polynomials, and the sum of the primitive roots modulo a prime.
At $s=1$ a different ingredient enters. A quantitative estimate for
$M(x)=\sum_{n\le x}\mu(n)$ gives ordinary convergence of three boundary
moments to $0,-1,-2\gamma$. We give the arithmetic proofs, an explicit
estimate obtained from ProveIt's finite zero-free-strip argument, and the
translation of the limits into primitive recursive precision statements.
The formalization section locates existing embedded derivations and describes the
steps to formalize the additional results in Lean.
\vspace{3pt}
\begingroup\small\tableofcontents\endgroup
\clearpage
\section{Arithmetic language and constructive factorization}\label{sec:arith}

Throughout, $[P]$ denotes the bit of a decidable predicate. Every divisor
is positive. The value of $\mu$ at zero is immaterial and may be set to
zero by an outer test. Write $\mathbf1(n)=1$ and $\varepsilon(n)=[n=1]$ on
positive inputs. A function in a theorem schema means a fixed
primitive recursive (PR) definition with its numerical parameters free,
or a supplied finite table covering the computation.

\paragraph{The object theory.}
We use first-order intuitionistic arithmetic with defining equations
for PR functions, the successor axioms, and induction on quantifier-free
formulas. An integer can be represented by $(a,b)$ with value $a-b$.
Equality to $(c,d)$ is $a+d=b+c$. Multiplication uses
$(ac+bd,ad+bc)$. Rationals have a positive denominator and equality by
cross multiplication. Finite lists can be built by iterated pairing,
with length carried separately, so every access is a bounded recursion.
Polynomial coefficients, residue classes, tables of finite operations,
and tuples of rationals have codes of this kind.
For a decidable test $R(i,\mathbf u)$ and a PR bound $B$,
$\prod_{i<B}[R(i,\mathbf u)]=1$ expresses that every test passes.
The sum of the bits has positive sign if a witness exists, and the first
witness is found by a scan of $B$ entries. A proof by strong induction
on a decidable matrix uses ordinary induction on the product of all its
preceding bits. Finite sums below are reindexed by explicit maps or by
partitioning a bounded rectangle into its fibers.

The real notation abbreviates arithmetic statements about rational
approximants. For example, a tail estimate means a universal statement
in the precision and the two truncation indices, after clearing positive
denominators and adding the evaluation tolerance. Complex
inputs may be Gaussian rationals or fixed PR Cauchy names with a supplied
modulus, size bound and the indicated domain margins. Proofs about such
names are schemata in the defining algorithms. A finite input version
takes the needed samples and tests the finitely many error bounds in the
proof. No quantification over infinite sequences is part of iPRA.

\paragraph{Division and prime factors.}
For $n>0$, define $p\mid n$ by a zero remainder for $p>0$.
Primality of $p\ge2$ is the absence of divisors between $2$ and $p-1$.
Define $v_p(n)$ for $p\ge2$ by the largest $j\le n$ such that $p^j\mid n$.
The search is bounded because $2^j\ge j+1$. Every $n>1$ has a least
divisor $p>1$, obtained by scanning through $n$. A proper factor of
$p$ would contradict leastness, so $p$ is prime. Repeatedly divide the
current remainder by its least prime, stopping at $1$ and then padding
with idle steps. A nonidle step strictly decreases the remainder, hence
$n$ steps suffice. The product of extracted primes times the remainder
is always $n$.

For uniqueness, the Euclidean divisions of $a,b$ strictly decrease the
nonzero remainder. Back substitution through this bounded list gives
integers $u,v$ with $ua+vb=\gcd(a,b)$. Initially represent $a,b$ by
$(1,0),(0,1)$ and subtract quotient times the preceding representation
at each division. Thus $p\nmid a$ implies $ua+vp=1$ for prime $p$, so
$p\mid ab$ forces $p\mid b$. Apply this lemma successively along any
second list representing $n$, cancel the matching prime, and induct
on list length. This proves uniqueness up to permutation, additivity
$v_p(ab)=v_p(a)+v_p(b)$, and the bijection between divisors of $n$
and exponent vectors $0\le e_p\le v_p(n)$. Both directions of this
bijection are bounded computations (take valuations, or multiply prime
powers). These arguments can also start from ProveIt's valuation and
Bézout derivations.

\section{Cancellation, inversion, and counting}\label{sec:core}
Set
\[
 \omega(n)=\sum_{\substack{p\le n\\p\text{ prime}}}[v_p(n)>0],\quad
 \Omega(n)=\sum_{\substack{p\le n\\p\text{ prime}}}v_p(n),\quad
 S(n)=\prod_{\substack{p\le n\\p\text{ prime}}}[v_p(n)\le1].
\]
Define $\lambda(n)=(-1)^{\Omega(n)}$ and $\mu(n)=S(n)(-1)^{\omega(n)}$.
Thus $\mu(1)=1$, $\mu(p)=-1$, $\mu(p^a)=0$ for $a\ge2$.
A square divisor $t^2$ with $t>1$ supplies a prime $p\mid t$. Conversely
$p^2\mid n$ is already such a divisor. This proves
\[
 |\mu(n)|=\mu(n)^2=S(n),\qquad
 \mu(n)=[\omega(n)=\Omega(n)]\lambda(n).
\]
Indeed every positive exponent contributes one to $\omega$, and the
excess in $\Omega$ is a sum of nonnegative integers. If $a,b$ are
coprime their prime supports are disjoint, so both the squarefree
indicator and the sign factor multiply. If $\gcd(a,b)>1$, its least
prime divisor has square dividing $ab$, giving $\mu(ab)=0$.
All these alternatives are decidable. For instance $30=2\cdot3\cdot5$
is the least composite with $\mu=-1$. A squarefree composite with an odd
number of prime divisors needs at least three, and its three smallest
are at least $2,3,5$. The same argument gives $2310=2\cdot3\cdot5\cdot7\cdot11$
as the first integer with five distinct prime factors.

\begin{theorem}[Divisor cancellation]
For $n\ge1$,
\[
 \sum_{d\mid n}\mu(d)=\varepsilon(n),\qquad
 \sum_{d\mid n}|\mu(d)|=2^{\omega(n)}.
\]
If $n>1$, exactly $2^{\omega(n)-1}$ divisors have each nonzero sign.
\end{theorem}
\begin{proof}
List the distinct prime divisors of $n$ in increasing order. Its squarefree
divisors are in bijection with bit masks of length $\omega(n)$ by the
product of the selected primes. For $n>1$ flip the first bit. This is
a fixed-point-free involution reversing the parity of the number of
selected bits. Exactly one member of each pair has its first bit clear.
The other bits are unrestricted. For $n=1$ there is one empty mask.
Grouping the masks by their number of selected bits gives also
$\sum_j(-1)^j\binom rj=[r=0]$.
\end{proof}

Define the Dirichlet convolution
$(f*g)(n)=\sum_{d\mid n}f(d)g(n/d)$. It is commutative for scalar functions by
$d\leftrightarrow n/d$, associative because both bracketings enumerate
the same ordered triples with product $n$, and has unit $\varepsilon$.
For coprime $u,v$, the divisors of $uv$ are exactly $ab$ with $a\mid u$,
$b\mid v$. Factoring the two finite sums shows that convolution preserves
multiplicativity. The theorem says $\mu*\mathbf1=\varepsilon$.
It uniquely specifies $\mu$. Its value at $1$ must be $1$, and at
$n>1$ it is minus the sum over proper divisors. This triangular recursion
also computes $\mu$ without factoring, by building a table through $n$.
\begin{theorem}[Two-sided inversion]
For fixed integer-valued PR functions, $F(n)=\sum_{d\mid n}f(d)$ for
all $n\ge1$ if and only if
$f(n)=\sum_{d\mid n}\mu(d)F(n/d)$ for all $n\ge1$.
The same holds up to any fixed bound and in a coded abelian group.
\end{theorem}
\begin{proof}
Convolve the forward equality with $\mu$, or the reverse one with
$\mathbf1$, and use associativity and the unit. On a finite interval
all arguments on the right are at most the point being evaluated.
For group-valued $f$, expand integer multiples by repeated addition
and inverse and use the same finite reindexing.
\end{proof}
In particular $\tau=\mathbf1*\mathbf1$ counts the divisors, and
$\mu*\tau=\mathbf1$. Also $\sum_{d\mid n}\varphi(d)=n$. The integers
$1\le k\le n$ with $n/\gcd(k,n)=d$ correspond to the residues coprime
to $d$ by dividing $k$ by $n/d$. Inversion gives
$\varphi(n)=\sum_{d\mid n}\mu(d)n/d
 =n\prod_{p\mid n}(1-1/p)$, by expansion over the prime masks.
Clearing denominators gives the integer formulation.

For $N\ge1$, group the pairs $dk\le N$ by their product, or by $d$, to get
\[
 1=\sum_{d\le N}\mu(d)\lfloor N/d\rfloor
  =\sum_{k\le N}M(\lfloor N/k\rfloor),\quad
 M(0)=0,\quad M(N)=1-\sum_{k=2}^N M(\lfloor N/k\rfloor).
\]
The prime masks in $\sum_{d^2\mid n}\mu(d)$ give one factor $1-1$ for
every $p^2\mid n$ and an empty product otherwise. Consequently
\[
  S(n)=\sum_{d^2\mid n}\mu(d),\qquad
  Q(N):=\sum_{n\le N}S(n)
     =\sum_{d\le\lfloor\sqrt N\rfloor}\mu(d)\lfloor N/d^2\rfloor.
\]
For $D=\lfloor\sqrt N\rfloor\ge1$ the floor errors total at most $D$,
and $\sum_{d>D}d^{-2}\le1/D$ by comparison with
$1/(d(d-1))$. Thus $|Q(N)/N-1/\zeta(2)|\le D/N+1/D\le2/D$,
using the Dirichlet identity proved below. This gives squarefree density
$6/\pi^2$ with a modulus. The counts with signs $+$ and $-$ are
$(Q(N)\pm M(N))/2$. Their common density is $3/\pi^2$ by the mean
estimate below.

\paragraph{Occupation-number interpretation.}
A vector of prime occupation numbers with finite support maps to
$n=\prod_p p^{a_p}$ and has energy $\sum a_p\log p=\log n$.
The inverse is the valuation vector. Restricting $a_p$ to $0,1$ selects
the squarefree integers. Their parity sign is $(-1)^{\sum a_p}$.
Assigning zero to excluded integers gives exactly $\mu$.
For a cutoff $n\le N$ this is an explicit bijection of finite tables.
Removing the cutoff under the absolutely convergent weights $n^{-s}$
gives the partition sums $\zeta(s)$, $\zeta(s)/\zeta(2s)$, and
$1/\zeta(s)$ in the unrestricted, restricted, and signed cases.

\section{Finite Fourier sums and Gauss's congruence}
\label{sec:fourier}

\paragraph{A weighted hyperbola identity.}
Fix $N\ge1$ and a table $w(1),\ldots,w(N)$ of integers. Partition the
pairs of positive integers $j,k$ with $jk\le N$ by their product. Each
pair in the part with product $u$ has the unique form $(u/k,k)$ with
$k\mid u$. Consequently
\[
   \sum_{jk\le N} w(jk)\mu(k)
    =\sum_{u=1}^N w(u)\sum_{k\mid u}\mu(k)=w(1).
\]
The constant table gives the floor-sum identity. The table
$w(u)=0,1,0,-1$ on residues $0,1,2,3\pmod4$ gives
$\sum_{jk\le N}\sin(\pi jk/2)\mu(k)=1$. These four integer values
supply the exact arithmetic interpretation of the sine. Each follows also
from the imaginary part of $i^u$.
Writing $W(k)=\sum_{1\le j\le N/k}w(jk)$ and $W(N+1)=0$, one more
finite rearrangement gives $\sum_{k=1}^N M(k)(W(k)-W(k+1))=w(1)$.
Indeed the coefficient of $\mu(d)$ after expansion telescopes to $W(d)$.
The same four-entry table therefore gives a Mertens version of the sine sum.


\paragraph{A polynomial certificate for the sum of primitive roots.}
Use rational polynomials stored as finite coefficient lists. For $n\ge1$
put $A_n=X^n-1$ and form the monic polynomials
\[
 D_n=\gcd\!\left(A_n,\prod_{e\mid n,\ e<n}(X^e-1)\right),
 \qquad H_n=A_n/D_n .
\]
An empty product is $1$. The gcd uses Euclidean division, which takes at
most one more than the degree of the smaller nonzero operand in division steps once the inputs are ordered by degree.
Since $A_n'=nX^{n-1}$ is coprime to $A_n$, the polynomial $A_n$ has no
nonconstant square divisor, since a repeated factor divides the derivative
by the product rule. For every proper divisor $e$ of $n$, the polynomial
$X^e-1$ divides both gcd operands and hence divides $D_n$. A common
factor of $H_n$ and $X^e-1$ would have its square dividing $H_nD_n=A_n$,
so their monic gcd is $1$. Here and below, one may verify coprimality by
the extended Euclidean algorithm. For $b>0$,
$\gcd(X^b-1,A_n)=X^{\gcd(b,n)}-1$ by repeated subtraction of
exponents, using
$X^a-1=X^{a-b}(X^b-1)+(X^{a-b}-1)$ for $a\ge b$.
In particular $\gcd(H_n,X^b-1)=1$ whenever $n\nmid b$.

Let $z$ denote the residue of $X$ modulo $H_n$, with equality tested by
polynomial division. Then $z^n=1$, and $z^b-1$ is invertible for every
positive $b$ with $n\nmid b$, by B\'ezout. For $h\ge0$, consider the
Ramanujan sum
\[
 C_n(h)=\sum_{0\le k<n,\ \gcd(k,n)=1}z^{hk}.
\]
Insert $[\gcd(k,n)=1]=\sum_{d\mid\gcd(k,n)}\mu(d)$. Writing $k=dj$ gives
an inner sum of $(z^{hd})^j$ for $0\le j<n/d$. If $n\mid hd$, each
summand equals $1$. Otherwise multiplication by the invertible $z^{hd}-1$
telescopes the sum to $z^{hn}-1=0$. Substitution of $e=n/d$ proves the
integer-valued formula
\[
       C_n(h)=\sum_{e\mid\gcd(n,h)}e\,\mu(n/e),
       \qquad C_n(1)=\mu(n).
\]
This is a uniform polynomial remainder identity, including $n=1$.
To evaluate at $\xi=\exp(2\pi i/n)$, each factor $\xi^e-1$ of the
proper-divisor product has modulus $2\sin(\pi\min(e,n-e)/n)>1/n$.
Indeed on $0\le t\le\pi/2<2$ the alternating sine series gives
$\sin t\ge t-t^3/6\ge t/3$. A product of $h$ such factors has modulus
at least $n^{-h}$. Writing the proper-divisor polynomial as $D_nV$,
the bound $|V(\xi)|\le 1+\sum_j|V_j|$ gives explicit apartness of
$D_n(\xi)$ from zero. Thus $A_n=D_nH_n$ forces $H_n(\xi)=0$.
Conversely every root of $H_n$ has exact order $n$ by the coprimalities
above. This evaluation gives the exponential formula, including the
convention $k=n$ in place of $k=0$.

\paragraph{The Farey formula.}
Let $\mathcal F_N$ consist of all reduced fractions in $[0,1]$ with
denominator at most $N\ge1$, both endpoints included. Sorting them has
no effect on a sum. Group the nonzero fractions by their denominator;
the group with denominator $d$ contains exactly the fractions $k/d$,
$1\le k\le d$, $\gcd(k,d)=1$. Its exponential sum is $\mu(d)$.
The extra term $0/1$ contributes $1$, yielding
\[
    \sum_{a\in\mathcal F_N}\exp(2\pi i a)=1+M(N).
\]
A single arithmetic certificate represents all summands in the residue
ring modulo $H_L$, $L=N!$. For denominator $d$, use $z^{L/d}$ in place
of the $d$th root. Its $b$th power equals $1$ if $d\mid b$, and
$z^{(L/d)b}-1$ is invertible otherwise. The same geometric-sum
calculation applies.

\paragraph{Generators modulo a prime.}
Let $p$ be prime and $m=p-1$. Multiplication by a nonzero residue $a$
permutes the $m$ nonzero residues of $\mathbb F_p$. Taking the product
and cancelling gives $a^m=1$. Its order is obtained by scanning
$1,\ldots,m$ for the first exponent giving $1$. Division with remainder
shows that this order divides every exponent giving $1$.
For each prime power $r^v$ exactly dividing $m$, the polynomial
$X^{m/r}-1$ has degree $m/r<m$; a nonzero field polynomial of degree $s$
has at most $s$ distinct roots, since division by $X-a$ removes a root
and reduces the degree by one. Scanning the nonzero residues therefore
returns an $a$ with $a^{m/r}\ne1$. The element $b_r=a^{m/r^v}$
has order $r^v$: its $r^v$th power is $1$, while its $r^{v-1}$st is
not. The product $g=\prod_r b_r$ has order $m$. Indeed if $g^u=1$,
raising to $m/r^v$ kills all the other factors and gives
$r^v\mid u(m/r^v)$, hence $r^v\mid u$. These divisibilities imply
$m\mid u$. For $m=1$ the empty product works.

The $m$ distinct powers $g^k$, $0\le k<m$, exhaust the nonzero
residues. Their orders are $m/\gcd(m,k)$, so the generators correspond
to the coprime exponents. Repeat the calculation of $C_m(1)$ with
$z=g$ in $\mathbb F_p$. Every nonzero difference is invertible by a
bounded field search. This proves Gauss's congruence
\[
  \sum_{\substack{1\le a<p\\ \operatorname{ord}_p(a)=p-1}} a
        \equiv \mu(p-1)\pmod p .
\]
In particular this yields $0$ when $p-1$ has a repeated prime factor,
and the parity sign otherwise.

\section{Counting irreducible polynomials over finite fields}
\label{sec:irreducibles}

Supply a finite field $F$ as a table of operations on $q\ge2$ labelled
elements, with distinguished $0,1$ and the field laws checked on the
table. A monic polynomial of degree $j$ is encoded by its $j$ free
coefficients, a base-$q$ word numbered from $0$ through $q^j-1$.
Let $I_q(j)$, $j\ge1$, count those with no monic divisor of degree $1$ through
$\lfloor j/2\rfloor$. The scan is exact since any proper factorization
has a smaller factor in this range and the leading coefficient can be
normalized. This defines a primitive recursive count.

\paragraph{Unique factorization suffices.}
A monic polynomial of positive degree either passes the irreducibility
test or splits into two monic factors of smaller positive degree. Repeating
the split gives irreducibles: the number of factors increases at each
split and cannot exceed the original degree. Over the field table,
long division cancels one leading term per step. Iterated division with
remainders, recording the two linear-combination coefficients, gives
B\'ezout's identity for a monic gcd. Thus if an irreducible $P$ divides
$AB$ and does not divide $A$, an identity $UP+VA=1$ multiplied by $B$
shows that $P$ divides $B$. It follows by stripping one irreducible at
a time and cancelling it that the multiset of monic irreducible factors
is unique. Every decision here is a comparison of finite codes.

\paragraph{The coefficient identity.}
Fix $n\ge1$ and work through degree $n$ in the ring
$\mathbb Z[T]/(T^{n+1})$. Enumerate all monic irreducibles $P$ of
degree $1$ through $n$ and put
\[
  B(T)=\sum_{j=0}^{n}q^jT^j
     =\prod_P (1-T^{\deg P})^{-1}.
\]
Each inverse is a finite geometric series in this ring. The coefficient
of $T^j$ in the product counts multisets of irreducibles of total degree
$j$, hence monic polynomials of degree $j$ by the preceding bijection.
The degree-zero monic polynomial is $1$. On this truncated ring the operator
$\Delta=T\,d/dT$ satisfies the product rule: it multiplies the $j$th
coefficient by $j$ and preserves the truncation. All factors in $B$ have
constant term one, so are units. Applying $\Delta B/B$ on both sides,
using $(1-qT)B=1$, gives
\[
  \frac{qT}{1-qT}
     =\sum_{d=1}^n I_q(d)\frac{dT^d}{1-T^d}.
\]
Comparison in degree $n$ gives $q^n=\sum_{d\mid n} d I_q(d)$.
Divisor inversion now yields
\[
             n I_q(n)=\sum_{d\mid n}\mu(d)q^{n/d}.
\]
The integer on the right is divisible by $n$ and nonnegative.
The proof uses only products of bounded length and coefficient arrays.

\paragraph{Constructing a field for every prime-power input.}
For completeness take $q=p$ prime first, with the residue field table
obtained by modular arithmetic. The coefficient count implies
$d I_p(d)\le p^d$ for $d\ge1$. For $r\ge2$ the contribution of all
proper divisors of $r$ is at most
$\sum_{d=1}^{\lfloor r/2\rfloor}p^d<p^r$; this follows by summing the
geometric progression and using $\lfloor r/2\rfloor+1\le r$.
Thus $I_p(r)>0$; for $r=1$ it equals $p$. A bounded scan of the $p^r$
monic degree-$r$ polynomials supplies an irreducible $P$. Remainders
of degree $<r$ modulo $P$ form the $p^r$-element field: B\'ezout
computes an inverse for every nonzero residue. The count just proved
can therefore be instantiated at each $q=p^r$ without asking for an
extension field or an irreducible polynomial as an oracle.

Conversely the size of every supplied finite field is a prime power. Among
the $q+1$ multiples $0,1,2\\cdot1,\\ldots,q\\cdot1$, two agree, so a
bounded scan gives a least positive $p$ with $p\\cdot1=0$. If $p=uv$
with $1<u,v<p$, the nonzero elements $u\\cdot1$ and $v\\cdot1$ have zero
product, contradicting the field laws. The multiples of $1$ form a copy
of $\\mathbb F_p$. Starting from the span $\\{0\\}$, scan the field for the
first element $z$ outside the current span $S$. Adjoining $z$ produces
the $p$ disjoint translates $S+az$ with $a\\in\\mathbb F_p$: equality
across two translates would express $z$ in $S$ by division by a nonzero
scalar. Each step multiplies the size by $p$, and after at most $q$ steps
the span is the whole field. Thus $q=p^r$ with $r\\ge1$.

\section{Incidence inversion and convolution powers}
\label{sec:incidence}
For a finite poset given by its decidable order table, fix an interval
$[a,b]$. Set $u(a,a)=1$ and, for $a<y\le b$, put
$u(a,y)=-\sum_{a\le x<y}u(a,x)$. To evaluate the recurrence,
order the elements by increasing cardinality of $[a,y]$, breaking ties
by their stored indices. Strict predecessors always have smaller
cardinality. Thus $\sum_{a\le x\le y}u(a,x)=[a=y]$, and each last
summand is determined by those preceding it. Equivalently, if the
order matrix in a linear extension is $Z=I+U$ on $m$ elements, the
strictly triangular matrix $U$ satisfies $U^m=0$. Multiplication shows
that $\sum_{j=0}^{m-1}(-U)^j$ is an inverse of $Z$ on both sides.
Its entries are the same $u$, with zero outside
comparable pairs. The proof applies to supplied finite intervals of a
locally finite poset whenever their order tables are given.
In the positive-integer divisibility order, the interval from $a$ to
$b$ with $a\mid b$ has exactly the elements $ad$ for $d\mid b/a$.
Hence $u(a,b)=\mu(b/a)$ by divisor cancellation and uniqueness.

For $k\ge0$, let $\mu_k$ be the $k$-fold Dirichlet convolution of
$\mu$, with $\mu_0(n)=[n=1]$. If $f,g$ are multiplicative then so is
$f*g$. Every divisor of a coprime product $ab$ corresponds uniquely to
a pair of divisors of $a,b$ by taking prime factors from the two disjoint
supports. The sum defining $(f*g)(ab)$ therefore separates into the two
sums at $a$ and $b$. At a prime power $p^v$, a nonzero term of the
$k$-fold convolution distributes the $v$ copies of $p$ among distinct
slots. Choosing the slots gives
\[
  \mu_k(p^v)=(-1)^v\binom{k}{v},\qquad
  \mu_k(n)=\prod_{p^v\parallel n}(-1)^v\binom{k}{v}.
\]
Here $p^v\parallel n$ in the product means $v>0$ is the exact exponent
and $\binom{k}{v}=0$ for $v>k$. For the polynomial extension, replace
$\binom{k}{v}$ by
$B_v(K)=K(K-1)\cdots(K-v+1)/v!$ in $\mathbb Q[K]$ with $B_0=1$.
The parameter identity $\mu_{K+L}=\mu_K*\mu_L$ follows locally from
$B_v(K+L)=\sum_{j=0}^v B_j(K)B_{v-j}(L)$: at nonnegative integer
arguments the equation counts $v$-element subsets of two disjoint
blocks. Each variable has degree at most $v$, so the finite grid
$0,\ldots,v$ in each variable proves equality of polynomials by the
root bound, applied first in one variable and then the other.
Prime-factor multiplicativity gives the general identity. These are
rational-coefficient array equalities in iPRA and specialize to complex
parameters under evaluation.

\section{Dirichlet series with a witnessed margin}
\label{sec:dirichlet-interior}

Take $h\ge1$ and a complex argument $s$ with
$\mathop{\rm Re}s\ge1+1/h$. A statement about this variable in iPRA is a
schema for a specified primitive-recursive (PR) rational approximation to
$s$, with bounds on its size and approximation error. For rational complex
$s$ the resulting statement is uniform in its numerator and denominator
codes. A disk or half-plane condition refers to a supplied bound.
Write $u_n=\exp(-s\log n)$, $Z(s)=\sum_{n\ge1}u_n$. Then
\begin{align}
 D(s):=\sum_{n\ge1}\mu(n)u_n
    &= Z(s)^{-1}=\prod_{p\ {\rm prime}}(1-u_p),\label{eq:mob-ds}\\
 Q(s):=\sum_{n\ge1}|\mu(n)|u_n
    &=Z(s)/Z(2s).\label{eq:mob-sqdir}
\end{align}
The notation $Z$ agrees here with the usual Dirichlet-series definition of
$\zeta$. There are explicit finite error bounds for each equality.

\paragraph{Evaluation of the elementary weights.}
For $n\ge1$ put $v=(n-1)/(n+1)$. Use the series
$2\sum_{j\ge0}v^{2j+1}/(2j+1)$ for $\log n$. After the terms
$j<L$ its remainder is at most $2v^{2L+1}/(1-v^2)$.
If $0\le r\le1-1/b$ and $b\ge2$, then $r^b\le1/2$ by
$(1+1/b)^b\ge2$ and $1-1/b\le(1+1/b)^{-1}$.
Thus $b=\max(2,n+1)$ supplies a PR cutoff.
For $|z|\le B$, $B$ a positive integer, the exponential Taylor terms
from index $2B$ on have successive absolute ratio at most $1/2$,
and the term at $2B$ is at most $(1+B)^{2B}$.
For example, the tail from $2B+L$ on is bounded by
$2(1+B)^{2B}2^{-L}$. These estimates also justify replacing $s$ and
$\log n$ by sufficiently accurate rationals. On a ball of radius $B$
the exponential is Lipschitz with bound $3^B$ by its power series
and $\sum 1/j!\le3$.

For clarity, the algebra behind the weights is elementary too.
The binomial theorem in the product of the exponential series gives
$\exp(a+b)=\exp(a)\exp(b)$. As a formal identity in $v$,
\[
 \exp\!\left(2\sum_{j\ge0}\frac{v^{2j+1}}{2j+1}\right)
       =\frac{1+v}{1-v}.
\]
Indeed both sides have constant term $1$ and solve
$(1-v^2)y'=2y$, which determines the coefficients successively.
At $0\le v<1$ the nonnegative terms can be collected by degree:
truncate the inner sum using its geometric remainder, truncate the outer
exponential using the factorial remainder above, then take a degree cutoff
at the product of the two polynomial degrees. This proves
$\exp(\log n)=n$. On real arguments $\exp$ is strictly increasing with
a lower increment bound on each bounded interval, by
$\exp(t)\ge1+t$ for $t\ge0$ and the addition law. These facts give
$u_{ab}=u_a u_b$. Conjugation and the addition law give
$|\exp(z)|=\exp(\mathop{\rm Re}z)$.

\paragraph{One estimate for all the rearrangements.}
In $2^j\le n<2^{j+1}$ the sum of $|u_n|$ is at most
$2^{-\lfloor j/h\rfloor}$. Indeed there are $2^j$ terms and
$n^{1/h}\ge2^{\lfloor j/h\rfloor}$, as seen on taking the $h$th power.
Set $B_h=2h$. For $T=2^{hL}$ every finite sub-tail on $n\ge T$
has absolute mass at most
\[
 E_h(L)=2h\,2^{-L},\qquad \sum_{n\ge1}|u_n|\le B_h .
\]
For a finite upper cutoff each bound follows by grouping $h$ successive
blocks and summing a finite geometric progression. The same bounds hold
with coefficients $\mu$, $|\mu|$, or any coefficients of absolute value
at most $1$.

In the product $D(s)Z(s)$ restrict first to $ab\le T^2$.
The coefficient of $u_m$ is $\sum_{a\mid m}\mu(a)=[m=1]$.
Every discarded pair has $a>T$ or $b>T$; its total absolute contribution,
summed over both strips, is at most $2B_h E_h(L)$. This proves $D(s)Z(s)=1$.
In particular $|Z(s)|\ge1/B_h$, so division has a rational apartness
bound. Expand the finite product over $p\le T$ by choosing one summand
from each factor. Unique factorization identifies its terms with
$\mu(m)u_m$ over squarefree $m$ whose prime factors are all at most $T$.
Every missing term has $m>T$. Hence the product approximates $D(s)$ with
error at most $E_h(L)$. Grouping subsets by size also gives the
alternating sums over increasing prime tuples printed by MathWorld.

Finally, each $m$ has exactly one decomposition $m=a b^2$ with $a$
squarefree: split each prime exponent into its parity and twice its
integer half. In $Q(s)Z(2s)$ this gives coefficient $1$ at $u_m$.
Truncate at $ab^2\le T^3$; the discarded pairs lie in $a>T$ or $b>T$,
and $|u_{b^2}|\le |u_b|$. The same error bound proves
$Q(s)Z(2s)=Z(s)$. Applying the reciprocal identity at $2s$ permits
the division in \eqref{eq:mob-sqdir}.

\section{Lambert series and the exponential product}
\label{sec:mob-lambert}

These identities hold coefficient by coefficient in the truncated
polynomial rings $\mathbb Z[X]/(X^{H+1})$ or
$\mathbb Q[X]/(X^{H+1})$ as appropriate. They also hold as effective
complex identities on every disk $|q|\le \rho=1-1/b$, $b\ge2$:
\begin{align}
 \sum_{d\ge1}\frac{\mu(d)q^d}{1-q^d}&=q,\label{eq:mob-lam1}\\
 \sum_{d\ge1}\frac{\mu(pd)q^d}{q^d-1}
    &=\sum_{j\ge0}q^{p^j}\qquad(p\text{ prime}),\label{eq:mob-lam2}\\
 \prod_{d\ge1}(1-q^d)^{\mu(d)/d}&=\exp(-q).\label{eq:mob-logprod}
\end{align}
In the last identity the factor means
$\exp((\mu(d)/d)\log(1-q^d))$ with the logarithm given by its series
at $1$. This specifies its complex branch.

In \eqref{eq:mob-lam1} expand the geometric sum. The coefficient at
degree $m$ is $\sum_{d\mid m}\mu(d)=[m=1]$.
For \eqref{eq:mob-lam2} it is $-\sum_{d\mid m}\mu(pd)$.
Writing $m=p^a v$, $p\nmid v$, the summands with $p\mid d$ vanish,
while $-\mu(pd)=\mu(d)$ for $p\nmid d$. The answer is
$\sum_{d\mid v}\mu(d)=[v=1]$, including $m=1=p^0$.
For \eqref{eq:mob-logprod} add the logarithms:
\[
 -\sum_{d,j\ge1}\frac{\mu(d)}{dj}q^{dj}
 =-\sum_{m\ge1}\frac{q^m}{m}\sum_{d\mid m}\mu(d)=-q.
\]
Modulo $X^{H+1}$ all displayed sums are finite. For the product use
the truncated exponential $\sum_{i=0}^H y^i/i!$ at an argument
with zero constant term. Its addition law is the binomial theorem in
degrees at most $H$.

For the analytic readings the double sums are dominated by
$\sum_{d,j\ge1}\rho^{dj}$. The strip $d>K$ has mass at most
$\rho^{K+1}/(1-\rho)^2$ by summing first over $j$, and the strip
$j>K$ has the same bound. Thus cutting at total degree $dj\le K^2$
leaves error at most $2b^2\rho^{K+1}$. The outer-row tail alone is
at most $b^2\rho^{K+1}$. For target $2^{-t}$ one may take
$K=b(t+2b+2)$, since $\rho^b\le1/2$.
The lacunary sum uses the same degree cutoff and its remaining mass after
degree $K^2$ is at most $b\rho^{K^2+1}$.
In the logarithmic case each finite sum of logarithms has norm at most
$b^2$; exponentiating transfers the error with the factor $3^{b^2}$.
These estimates justify the exact rearrangements with only finite
geometric and factorial bounds.

\section{The three reciprocal-square masses}
\label{sec:mob-basel}

The values tabulated in MathWorld depend on two even zeta values. The
following finite polynomial proof supplies both and gives a PR rate.
We use angles in radians.

\paragraph{Effective trigonometry.}
Define sine and cosine by the imaginary and real parts of $\exp(it)$.
For $|v|<1$ put
$A(v)=\sum_{j\ge0}(-1)^jv^{2j+1}/(2j+1)$ and set
$\pi=4(A(1/2)+A(1/3))$. This is the usual arctangent normalization of
$\pi$. The alternating bounds give $2<\pi<4$.
The formal identity $\exp(2iA(v))=(1+iv)/(1-iv)$ follows by matching
the constant term and the coefficients of $(1+v^2)y'=2iy$.
At $|v|\le1/2$ absolute geometric and exponential majorants justify
evaluation as in Section~\ref{sec:dirichlet-interior}.
Multiplying the two rational values on the right for $v=1/2,1/3$
gives $i$. Thus $\exp(i\pi/2)=i$ and $\exp(i\pi)=-1$.

For $0\le x\le\pi/2<2$, the alternating series give
$x-x^3/6\le\sin x\le x$, in particular $\sin x\ge x/3$.
Also $\cos x=\sin(\pi/2-x)\ge0$ by the addition law, and
$|\sin x|,|\cos x|\le1$ by the complex norm identity.
We shall need $x\cos x\le\sin x$ on this interval. Indeed
$\sin x-x\cos x$ has successive signed terms
$(-1)^{j+1}2j x^{2j+1}/(2j+1)!$, $j\ge1$. Their absolute ratio is
$(j+1)x^2/[j(2j+3)(2j+2)]\le 2/5$ for $x\le2$.
Pairing successive terms proves the inequality. For $0<x<\pi/2$
we consequently have
\begin{equation}
 0\le x^{-2}-\cot^2 x\le1 .\label{eq:mob-cot-est}
\end{equation}
The upper bound is $x^{-2}\le1/\sin^2x=1+\cot^2x$, and the lower
bound follows from $x\cos x\le\sin x$.

\paragraph{A finite polynomial with known roots.}
Take $N=2m+1$ with $m\ge2$ and $x_k=k\pi/N$, $1\le k\le m$.
The polynomial $\mathop{\rm Im}(Y+i)^N$ (real indeterminate $Y$)
is an even polynomial in $Y$. Set $U=Y^2$ and divide by $N$ to obtain
the monic polynomial
\[
 P(U)=U^m-\frac{\binom N3}{N}U^{m-1}
           +\frac{\binom N5}{N}U^{m-2}+\cdots .
\]
For each $k$, $(\cot x_k+i)^N=(-1)^k/\sin^N x_k$, so
$r_k=\cot^2x_k$ is a root. There is explicit apartness between
different roots. For $x_j<x_k$,
$\cot x_j-\cot x_k=\sin(x_k-x_j)/(\sin x_j\sin x_k)
\ge \pi/(3N)$.
Further $\cot x_k\ge \cos x_k\ge(\pi/2-x_k)/3\ge\pi/(6N)$.
Thus differences of the $r_k$ have a positive rational lower bound.
Synthetic division successively extracts the $m$ linear factors:
evaluating a quotient at the next root gives zero because each previously
extracted factor there is apart from zero. This proves the coefficient
identities
\[
 \sum_k r_k=C,\qquad \sum_k r_k^2=C^2-2E,\qquad
 C=\frac{(N-1)(N-2)}6,\quad E=\frac{\binom N5}{N}.
\]

\paragraph{Passage to the sums.}
Set $t_k=(\pi/N)^2r_k$. By \eqref{eq:mob-cot-est},
$0\le k^{-2}-t_k\le (\pi/N)^2$. Their squares differ by at most
$2(\pi/N)^2$ since $t_k\le k^{-2}\le1$.
The tail of either $\sum k^{-2}$ or $\sum k^{-4}$ beyond $m$ is
at most $1/m$, using
$1/k^2\le1/(k-1)-1/k$.
Hence the respective sums differ from $(\pi/N)^2 C$ and
$(\pi/N)^4(C^2-2E)$ by at most $1/m+2m(\pi/N)^2$.
Now $C/N^2\to1/6$ and
$E/N^4=\prod_{j=1}^4(1-j/N)/120\to1/120$.
For an explicit check use
$0\le1-\prod_j(1-j/N)\le\sum_j j/N$ for $N\ge5$.
It gives $|C/N^2-1/6|\le1/(2N)$,
$|E/N^4-1/120|\le1/(12N)$ and
$|(C/N^2)^2-1/36|\le1/(6N)$.
All errors thus tend to zero with a PR modulus, e.g.\ taking
$m=128\cdot2^t$ makes each combined error less than $2^{-t}$.
We obtain $Z(2)=\pi^2/6$ and
$Z(4)=\pi^4(1/36-2/120)=\pi^4/90$.

At $s=2$ equations \eqref{eq:mob-ds}--\eqref{eq:mob-sqdir} now give
$D(2)=6/\pi^2$ and $Q(2)=15/\pi^2$. Since twice the indicator of
$\mu(n)=\pm1$ is $|\mu(n)|\pm\mu(n)$, the two separate masses are
\begin{equation}
 \sum_{\mu(n)=-1}\frac1{n^2}=\frac9{2\pi^2},\qquad
 \sum_{\mu(n)=1}\frac1{n^2}=\frac{21}{2\pi^2}.
\end{equation}
For all three positive series the rational partial sum through $m$
already has error at most $1/m$.

\paragraph{Lean implementation of the precision bounds.}
Represent a Gaussian rational by a pair of rational numerators with
positive denominators. For instance, the scalar cutoff
$\texttt{tailCut}(h,L)=2^{hL}$ and the disk cutoff
$\texttt{diskCut}(b,t)=b(t+2b+2)$ have type
$\mathbb N\to\mathbb N\to\mathbb N$. Define a finite series evaluation
from a list of coefficients, a cutoff and an input approximation.
The Dirichlet proof requires three finite lemmas: a bound for a dyadic
block, a bijection between divisor sums and hyperbola sums, and the
two-strip bound. An embedded derivation of the last lemma first scans the
rectangle by rows, includes all points with $a>T$, then includes only
$a\le T$ in the column strip. This uses a bounded sum of nonnegative
rational majorants. The supplied
\texttt{lean/FormalSeries.lean} verifies the ordinary Lambert coefficients
and the logarithmic coefficients of the product at every degree. It
represents $L$ by $l_n=n[X^n]L$ in positive degrees and represents $E$
by $e_n=n![X^n]E$. The formal equation $E'=L'E$ then reads
\[
 e_{n+1}=\sum_{k=0}^n n(n-1)\cdots(n-k+1)\,l_{k+1}e_{n-k},
\]
with an empty product at $k=0$. The file proves that $l_n=-[n=1]$
and that the normalized solution has $e_n=(-1)^n$.
For the remaining formal power series one may represent a
truncation by an array of $H+1$ rational codes. Cauchy multiplication
at index $k\le H$ is the bounded sum over $i\le k$. Formal differentiation
through degree $H$ reads an input through degree $H+1$.
Prove both coefficient ODE identities with that extra input coefficient.
The polynomial proof of the two zeta values uses coefficient descent in
synthetic division and the explicit lower bounds at the roots.
For real inequalities $A\le B$ with error radius $e$ on each input,
the target at a given precision is the rational comparison
$A_{\rm app}\le B_{\rm app}+2e$. Equality uses both inequalities.
The finite induction arguments can therefore be stated with a fixed precision
parameter and a bounded Boolean check for every earlier row, coefficient
or polynomial degree needed by the recursion.
Here is an error budget for inductions whose length depends on the
input. Generate the finite list of arithmetic operations and inference
steps for the chosen cutoffs; let $J$ be its length. Bound magnitudes
forward by replacing addition by the sum of bounds and multiplication
by their product. At each inverse or cancellation record the reciprocal
of the supplied lower bound. For example at the cotangent roots one may
use $|\sin x_k|\ge 2/(3N)$ and
$|r_j-r_k|\ge 4/(9N^2)$ for $j\ne k$.
Let $C\ge2$ be an integer exceeding all these bounds. With
$K=8(1+C)^4$, addition, multiplication, inversion on a ball of radius
less than $1/(2C)$ around an admissible denominator, and cancellation
against an apart factor all amplify incoming errors by at most $K$ per
step. Indeed a product error is at most $(2C+1)e$ and an inverse error
at most $2C^2e$ for input errors at most $e\le1/(2C)$.
Start every elementary numerical evaluation with error at most
$2^{-t}/(4CK^{J+1})$. Even with one new error of this size at each step,
the geometric sum bounds the final accumulated error by $2^{-t}$.
Power-series truncations are chosen by the preceding explicit tails.
All bounds, including the generation and length of the operation list,
use nested bounded loops or primitive recursion. This permits a single
quantifier-free induction on the list at the selected working precision,
including the synthetic-division argument with a variable number of roots.

\section{Cancellation at the boundary of the Dirichlet series}\label{sec:boundary}

Evaluation of the absolutely convergent series at $1+t$ settles none of
the three ordinary series at $t=0$ by itself. We prove a summatory estimate
strong enough to take that limit uniformly for two derivatives. We use
the finite analytic lemmas in ProveIt's logarithmic-integral PNT packet
\cite{strip}. We state the particular estimates, explain their proofs,
and give the change of integrand and the unsmoothing step in detail.

\subsection{A quantitative estimate for the signed sum}

\begin{theorem}\label{thm:strongM}
Put $D=2^{100}$. For an integer $x$ with $x\ge 2^{D\,8192^{10}}$,
compute the largest $k$ satisfying $2^{Dk^{10}}\le x$ by bounded search.
Then
\[
      |M(x)|\le x\,2^{-k/8}.
\]
In particular, $|M(x)|\le x/(\log x)^8$ for every integer
$x\ge X_0=2^{\,2^{300}}$. The first estimate has the exact arithmetical
form $2^k|M(x)|^8\le x^8$.
\end{theorem}

\paragraph{The strip and its numerical lower bound.}
Let $T=2^k$, $Q=k+2$, $\delta=2^{-88}Q^{-9}$ and $\alpha=1-\delta$.
On the rectangle $\alpha\le\sigma\le1+1/64$, $|t|\le T$,
the reciprocal $F(s)=1/\zeta(s)$ has a regular extension at $1$ and
\begin{equation}\label{eq:stripM}
                      |F(s)|\le 2^{71}Q^7 .
\end{equation}
Here $s=\sigma+it$. The numerical content of the strip lemma matters.
Euler summation, obtained by integration on each unit interval, gives
\[
 \zeta(s)=\sum_{n\le V}n^{-s}+\frac{V^{1-s}}{s-1}
          -s\int_V^\infty\{u\}u^{-s-1}\,du.
\]
The integral tail after $W$ is at most $|s|W^{-\sigma}/\sigma$.
Its first three derivatives are controlled by the same integrals with
up to three powers of $\log u$, evaluated by parts. With
$H=|t|+3$, $V=\lceil H^2\rceil$, and $|t|\ge1/64$, bounds (A1)
of \cite{strip} give $|\zeta|\le2^9\log H$ and
$|\zeta'|\le2^{15}(\log H)^2$ when $\sigma\ge1-1/(2\log H)$
and $\sigma\le2$. The finite sum has $V^{1-\sigma}\le4$ as its
weight bound, and the pole term has $|s-1|^{-1}\le64$.

For $\sigma>1$ multiply the Euler products with weights $3,4,1$ at
$t=0,t,2t$. The real logarithm of the product of their moduli has
summands proportional to $3+4\cos u+\cos 2u=2(1+\cos u)^2$,
all nonnegative. Thus
$\zeta(\sigma)^3|\zeta(\sigma+it)|^4|\zeta(\sigma+2it)|\ge1$.
It already suffices to certify the lower bound $1/2$ by finite
products. The tails are bounded by $V^{1-\sigma}/(\sigma-1)$.
Use $\zeta(1+\eta)\le2/\eta$ for $0<\eta\le1$ and the (A1)
upper bound $2^{10}Q$ at $2t$. At the rational point
$1+\delta+it$ the fourth-power inequality gives
$|\zeta|\ge2^{-4}\delta^{3/4}Q^{-1/4}=2^{-70}Q^{-7}$.
Moving left by $2\delta$ costs at most $2^{16}\delta Q^2$;
farther right the same direct lower bound only increases. This gives
\eqref{eq:stripM} for $|t|\ge1/64$.
For $|t|\le1/64$ write
\[
 h(s)=1-s\int_1^\infty\{u\}u^{-s-1}\,du,\qquad
 F(s)=\frac{s-1}{1+(s-1)h(s)} .
\]
There $|h|\le4$, $|s-1|\le1/32$ and the denominator has modulus at
least $7/8$. The formula includes $F(1)=0$ and agrees with the other
one on their overlap.

\paragraph{A signed smoothed sum.}
For rational $y\in[x,2x]$, put $b=\lceil\log_2 y\rceil$ (an integer
obtained by comparison of powers), $c=1+1/b$ and
$B(y)=\sum_{n\le y}(y-n)\mu(n)$. The scalar Perron kernel gives
\begin{equation}\label{eq:smoothM}
 B(y)=\frac1{2\pi i}\int_{c-i\infty}^{c+i\infty}
                F(s)\frac{y^{s+1}}{s(s+1)}\,ds .
\end{equation}
Indeed expand $F$ at $c$, using $\sum|\mu(n)|n^{-c}\le1+b$.
For $v=y/n$, the integral of $v^s/(s(s+1))$ on this line, with the
factor $1/(2\pi i)$, is $(1-1/v)_+$. For $v\ge1$ close to the
left past the poles $0,-1$ with residues $1,-1/v$. For $v<1$ close
to the right, with no poles. On the closing vertical line at $\pm U$
the integral is bounded by $1/(2(U-1))$, using $v^{\pm U}\le1$ on the
chosen side and integrating the quadratic denominator. The two horizontal
edges at height $W$ cost at most $(U+2)\max(1,v^2)/W^2$ up to a factor
of two. These bounds also handle $v=1$. First truncate the Dirichlet
series at $N$ using $3y^2bN^{-1/b}$ for its integrated error, then choose the kernel error per term
less than the required error divided by $1+Ny$. All cutoffs are PR.
The two residues can equivalently be evaluated by integrating the
explicit simple fractions around rational rectangles (lemma (E8) of
\cite{strip}).

Truncate \eqref{eq:smoothM} at height $T$. Since $y^{c+1}\le3y^2$,
the omitted tails cost at most $3(1+b)y^2/T$. The contour may be
moved to $\alpha$ without a residue, by \eqref{eq:stripM} and the
regular formula at $1$. The integral of $|s(s+1)|^{-1}$ along this
vertical segment is at most $\pi/\alpha\le2\pi$. On the horizontal
edges the denominator has absolute value at least $T^2$. With
$R=\log y$, we may therefore use the slightly larger common bounds
\begin{equation}\label{eq:smoothBound}
 \frac{|B(y)|}{y^2}
 \le 2^{91}Q^9e^{-\delta R}
       +2^{88}Q^9/T^2+54 R^2/T .
\end{equation}
The first term uses $y^{\alpha+1}=y^2e^{-\delta R}$.
These are the numerical majorants in (A6) of \cite{strip}, used here for
a different Dirichlet series.

For clarity, the contour move uses a strictly finite cancellation lemma.
On a rational rectangular grid the trapezoidal boundary sums of every
affine complex function vanish by $(u+v)(v-u)=v^2-u^2$. Approximation
by the first-order Taylor polynomial on each cell therefore gives a
total error bounded by $2JL^3/P+4LP\eta+2L^2\eta$, where $L\ge1$ bounds
the sum of the side lengths, $P$ is the number of equal divisions per
side, $J$ bounds the first two derivatives and the Taylor remainder,
and $\eta$ is the per-value error in the taxicab norm. Internal edges
cancel with opposite orientations. Take
$P=2^{v+8}(J+1)(L+1)^3$ and
$\eta\le 2^{-v-10}/((L+1)^2(P+1))$ for accuracy $2^{-v}$.
For this integrand one may take
$J=2^{512}(\lceil y\rceil+1)^8(T+4)^{40}$.
Away from $|t|<1/64$, Euler summation at $V=1$ bounds the first
two zeta derivatives and zeta itself by $2^{40}(T+3)$, and
\eqref{eq:stripM} bounds $|F|$ by $2^{71}(T+3)^7$.
Then $F'=-\zeta' F^2$ and
$F''=2(\zeta')^2F^3-\zeta''F^2$. In the central band differentiate
the regular quotient with denominator $\ge7/8$ using the integral bounds
for $h$. The derivatives through order two of $y^{s+1}/(s(s+1))$ are
bounded by $2^{10}y^6$ since $\sigma\ge1/2$.
The product rule gives the displayed $J$ with room for both norm
conversions. These derivative bounds apply on both sides of the rational
band boundaries with matching values. Rational evaluations use the
Euler-summation tail, log and exponential enclosures; rectangle edges
are integrated as finite sums with the preceding errors. Thus the move
does not require integration or existence of zeros as an object-theory axiom.

\paragraph{Choice of scales and removal of smoothing.}
We have $Q\le2k$, $2^{98}k^{10}\le R\le2^{112}k^{10}$,
so $\delta R\ge2k$. For $k\ge8192$, $k\le 2^{k/256}$.
Indeed with $u=\lfloor k/256\rfloor\ge32$, $256(u+1)\le2^u$ by
induction. It follows that
$2^{100}k^9\le2^{k/8}$ and $54\cdot2^{224}k^{20}\le2^{k/4}$:
after bounding $k$ by $2^{k/256}$ both inequalities are linear in $k$
in the exponent. Each of the first two terms of \eqref{eq:smoothBound}
is at most $2^{-15k/8}$ and the third at most $2^{-3k/4}$. Therefore
$|B(y)|\le y^2 2^{-k/2}\le4x^2 2^{-k/2}$.

Set $h=x\,2^{-\lfloor k/4\rfloor}$. For integer $x$,
\[
 \frac{B(x+h)-B(x)}h
 = M(x)+\sum_{x<n\le x+h}\mu(n)\frac{x+h-n}{h}.
\]
Each weight on the right is between $0$ and $1$ and the last sum has
at most $\lfloor h\rfloor$ terms. Its magnitude is at most $h$.
Consequently $|M(x)|\le8x^2 2^{-k/2}/h+h\le10x2^{-k/4}
\le x2^{-k/8}$. The difference estimate uses just $|\mu(n)|\le1$,
so it applies to these signed coefficients.

When $x\ge X_0$, $k\ge2^{20}$. Write $u=\lfloor\log_2 k\rfloor$.
The inequality $1024(u+1)\le2^u$ holds at $u=20$ and propagates by
doubling, hence $\log_2 k\le k/1024$. Thus
$8\log_2\log x\le896+80\log_2 k\le k/8$. This completes the proof
of Theorem~\ref{thm:strongM}.

\subsection{Three convergent boundary moments}

\begin{theorem}
With partial sums taken in increasing integer order,
\[
 \sum_{n\ge1}\frac{\mu(n)}n=0,\qquad
 \sum_{n\ge1}\frac{\mu(n)\log n}n=-1,\qquad
 \sum_{n\ge1}\frac{\mu(n)(\log n)^2}n=-2\gamma.
\]
Here $\gamma=\lim_{N\to\infty}(H_N-\log N)$, $H_N=\sum_{j=1}^N 1/j$.
\end{theorem}
\begin{proof}
For $j=0,1,2$ and $0\le t\le1$, use $w(u)=(\log u)^j/u^{1+t}$,
taking the zeroth power to be one at $u=1$.
For $u\ge N\ge X_0$ we have
$|M(\lfloor u\rfloor)|\le256u/(\log u)^8$, since
$\lfloor u\rfloor\ge\sqrt u$.
Also $|w'(u)|\le4(\log u)^j/u^2$.
For integers $V\ge N$, summation by parts gives
\[
 \sum_{N<n\le V}\mu(n)w(n)=M(V)w(V)-M(N)w(N)
                    -\int_N^V M(\lfloor u\rfloor)w'(u)\,du .
\]
Each unit-interval integral simply means the finite difference of $w$
times the corresponding constant. Bounding the endpoints and integrating
$1024/(u(\log u)^6)$ gives the uniform estimate
\begin{equation}\label{eq:momentTail}
 \left|\sum_{N<n\le V}\mu(n)w(n)\right|\le
                       \frac{2048}{(\log N)^5}.
\end{equation}
In particular $N(v)=2^{\,2^{300}+2^{v+16}}$ makes this $<2^{-v-2}$.
This modulus is chosen for simplicity. All constants are integers and
the huge cutoff is used symbolically.

For $t>0$, the moment with index $j$ equals $(-1)^j F^{(j)}(1+t)$.
One can differentiate the absolutely convergent series twice on
$t\ge a>0$ by majorizing the log-weighted derivative tails with
$\int_N^\infty(\log(2u))^r u^{-1-a}\,du$, $r\le3$, evaluated by parts.
By \eqref{eq:momentTail} the moments extend continuously from the
right to the ordinary sums at $t=0$. Explicitly, first cut off at $N(v)$;
the head changes by at most $tN(v)(1+\log N(v))^3$ by
$|1-e^{-z}|\le z$ for $z\ge0$; then allow for both tails. This gives
a PR choice of positive rational $t$ at any desired precision.

At $s=1+t$ the regular formula is $F(1+t)=t/(1+t h(1+t))$.
Euler summation at $s=1$ yields
$H_N-\log N=1-\int_1^N\{u\}u^{-2}du$, so
$|H_N-\log N-h(1)|\le1/N$ and $h(1)=\gamma$.
Differentiating the regular quotient at zero gives successively
$F(1)=0$, $F'(1)=1$, $F''(1)=-2h(1)$. The first three derivatives
of $h$ near $1$ are bounded by its differentiated integrals, so the
right-hand limits of these quotient derivatives have PR error bounds.
Multiplying by $(-1)^j$ proves all three evaluations.
\end{proof}

For an explicit arithmetic output of the last theorem, take $n\ge N(v+3)$,
use dyadic logarithm enclosures of radius
$2^{-v-12}/(n+1)^3$ in the head through $n$, and approximate $\gamma$
by $H_{2^{v+7}}-\log(2^{v+7})$ to the same precision in its logarithm.
For $j\le2$ every logarithm in the head has size at most $n$, so the
change from raising to the $j$th power and summing is at most
$n(2n+1)2^{-v-12}/(n+1)^3$. Adding the tail and the error of the target
gives a rational distance $<2^{-v}$. This last assertion, with a
cross-multiplied integer matrix, is a first-order formulation of
ordinary convergence to each displayed value.

\subsection{A conditional link with the position of zeta zeros}

There is a precise constructive version of the article's observation
that the growth of $M$ controls zeros. Suppose one is given rational
$0\le\theta<1$ and an integer $C\ge1$ such that
$|M(n)|\le C n^\theta$ for every $n\ge1$. Denominators in the exponent
can be eliminated by taking integer powers. Then for rational complex
$s=\sigma+i\tau$, $\sigma>\theta$ and $s\ne1$,
\[
       |\zeta(s)|\ \ge\ \frac{\sigma-\theta}{C|s|}.
\]
In particular, a family of such bounds for every rational
$\theta>1/2$ with $\theta<1$ would exclude zeta zeros to the right
of the critical line. This conclusion is conditional on the stated
power bounds; Theorem~\ref{thm:strongM} supplies logarithmic savings.

To prove it, put $D_M(s)=s\int_1^\infty M(\lfloor u\rfloor)u^{-s-1}du$.
Its tail after $N$ is at most $|s|CN^{-(\sigma-\theta)}/(\sigma-\theta)$.
Finite partial summation identifies it with $\sum\mu(n)n^{-s}$, and
the integral bound gives $|D_M(s)|\le C|s|/(\sigma-\theta)$.
Let $W(s)=1+(s-1)h(s)$ with $h$ as above; thus $W=(s-1)\zeta$.
We verify that $G=W D_M-(s-1)$ vanishes throughout $\sigma>\theta$.

It suffices to explain the propagation from $2+i\tau$ to
$\theta<\sigma\le2$. Choose rational
$0<\delta\le\min(1/2,\sigma-\theta)$, put $r=\delta/2$, and move
horizontally in rational steps of size at most $r/4$.
At each center $z$, expand $u^{-z-v-1}$ in $v$ in both integrals.
The coefficient in degree $j$ contains $( -\log u)^j/j!$.
The sum of the absolute integral majorants at $|v|=r$ is bounded
using $u^r$, because the exponential series sums to it. Equivalently
one can use
$\int_1^\infty(\log u)^j u^{-1-a}du=j!/a^{j+1}$ for $a>0$.
Thus for the coefficients $a_j$ of $G(z+v)$ one has
$\sum_j |a_j| r^j\le B$, where, for example,
$H=6+\lceil|\tau|\rceil$ and
$B=\lceil16(C+1)(H+1)^4/\delta^2\rceil$ suffice.
Indeed the majorants for $D_M$ and $W$ there are respectively
$2HC/\delta$ and $1+H(1+2H/\delta)$, and products multiply majorants.

At $2+i\tau$ all coefficients vanish: expand the absolutely convergent
product $\zeta D_M$ termwise on its disk and group by the integer product
using divisor cancellation. At a step $v$ the new coefficient in
degree $\ell$, scaled by $r^\ell$, is the sum of
$\binom j\ell a_j r^j(v/r)^{j-\ell}$ over $j\ge\ell$.
The omitted part past $J\ge\ell$ is bounded by
$B4^\ell2^{-J}$ using $\binom j\ell\le2^j$. This makes propagation
of zero an arithmetic argument even for a variable number of centers.
Starting with the desired final degree zero and error $e$, work backwards:
for each step requiring coefficients through degree $l$ with error $e$,
take $J\ge l$ such that $B4^l2^{-J}\le e/2$ and ask for the preceding
scaled coefficients through $J$ to error $e/(2\cdot4^{l+1})$.
Indeed the sum of the absolute multipliers for one such coefficient
is at most $4^l\sum_{j\ge0}2^{-j}\le 2\cdot4^l$.
Choose $J$ by the binary length of the rational threshold, so the entire
schedule is PR. Now a forward induction on the bounded tables proves
$G(s)=0$ at every requested precision. For $s\ne1$ the rational
distance from $1$ supplies a positive bound for division, yielding
$\zeta(s)D_M(s)=1$ and the claimed inequality.

\section{What the mean-zero assertion says about primes}\label{sec:pnt}
The zero mean also has a two-way elementary connection to the prime
number theorem. The following proof transports an input modulus, so
the equivalence has content independently of Theorem~\ref{thm:strongM}.
Write $\psi(x)=\sum_{p^a\le x,\ a\ge1}\log p$ and
$\Lambda(n)=[n=p^a]\log p$ with the prime recovered by bounded
factorization and value zero if there is no such prime.
Write $m(x)=\sum_{d\le x}\mu(d)/d$ and take $x$ to be a positive
integer unless a real argument is indicated.

\begin{theorem}
The following limits imply one another with PR transformations of any
supplied PR modulus:
\[
  \frac{\psi(x)}x\longrightarrow1,\qquad
  \frac{M(x)}x\longrightarrow0,\qquad m(x)\longrightarrow0.
\]
The first limit is equivalent, in the same sense, to
$\pi(x)\log x/x\longrightarrow1$, where $\pi(x)$ counts primes.
\end{theorem}

\paragraph{From primes to the mean.}
The operation $Df(n)=f(n)\log n$ obeys the product rule for
convolution by $\log(ab)=\log a+\log b$.
Unique factorization gives $D\mathbf1=\mathbf1*\Lambda$.
Apply $D$ to $\mu*\mathbf1=\varepsilon$ and convolve with $\mu$.
It follows that $D\mu=-\mu*\Lambda$. On summing through $x$ we get
\[
 \sum_{n\le x}\mu(n)\log n=-\sum_{d\le x}\mu(d)\psi(x/d).
\]
By the floor identity, $xm(x)=1+\sum_{d\le x}\mu(d)\{x/d\}$,
so $|m(x)|\le1+1/x$. Also
$\sum_{n\le x}\log(x/n)\le\int_0^x\log(x/t)\,dt=x$
by monotonicity (the first term can instead be kept separately).
Given a rational $\eta\in(0,1)$, take $A\ge2$ from the PNT modulus
so that $|\psi(y)-y|\le\eta y$ for $y\ge A$. For real $y$ use
$\psi(y)=\psi(\lfloor y\rfloor)$ and increase the integer cutoff
beyond $2/\eta$, using the initial modulus at $\eta/2$.
For $1\le y<A$ the crude bound is $|\psi(y)-y|\le A^2+A$, by
$\Lambda(n)\le\log n$. Thus
\[
 \frac{|M(x)|}{x}\le
 \frac{A^2+A+3+\eta(1+\log x)}{\log x}\qquad(x\ge2).
\]
Indeed the main term after replacing $\psi(x/d)$ by $x/d$ is at most
$x|m(x)|$, and the remaining large-argument errors cost at most
$\eta x\sum_{d\le x}1/d$. For accuracy $1/k$ take
$\eta=1/(4k)$ and, for example, $x\ge2^{8k(A^2+A+5)}$.

\paragraph{A finite summation-by-parts lemma.}
Suppose $M(x)/x\to0$ with a modulus. For each of the three functions
of a rational $y\ge1$ displayed below one has
\[
 \frac1x\sum_{d\le x}\mu(d)g(x/d)\longrightarrow0:
 \quad
 \begin{cases}
  g_1(y)=\{y\},\\
  g_2(y)=y(H_{\lfloor y\rfloor}-\log y-\gamma),\\
  g_3(y)=\log(\lfloor y\rfloor !)-y\log y+y .
 \end{cases}
\]
Here are bounds that prove the lemma without a general equidistribution
or integration theorem. Each $|g(y)|\le8\sqrt y$. For $g_2$ use
$|H_j-\log j-\gamma|\le1/j$ and
$\log(y/j)\le1/j$ for $j=\lfloor y\rfloor$.
For $g_3$ the increasing function $\log t$ bounds $\log(j!)$ between
$\int_1^j\log t\,dt$ and $\int_1^j\log t\,dt+\log j$.
Replacing the upper integration endpoint $j$ by $y$ costs at most
$\log y$. Thus $|g_3(y)|\le 1+2\log y\le8\sqrt y$.
On $[1,A]$, $A$ integral, the two endpoint magnitudes plus the total
variation of each $g$ are bounded by $V_A=50(A+1)^4$.
For a direct check, on the open unit intervals the three derivatives are
$1$, $H_j-\log y-\gamma-1$, and $-\log y$. Their absolute values are
at most $3(A+1)$. The jumps at integers have magnitudes at most
$1,1,\log A$, respectively. The same bounds hold for finite differences
by integrating the explicit derivatives on each intervening piece.

The terms $d\le x/A$ in the original sum contribute at most
$16x/\sqrt A$, using $\sum_{d\le u}d^{-1/2}\le2\sqrt u$.
For the other terms, partial summation with the bounded variation on
$[1,A]$ costs at most $\eta x V_A$ if $|M(t)|\le\eta t$ for all
integers $t\ge\lfloor x/A\rfloor$. Thus for a target $\epsilon>0$
choose $A\ge(32/\epsilon)^2$, then $\eta=\epsilon/(2V_A)$, and take
$x\ge A(U+1)$ where $U$ is the input modulus at $\eta$.
All variation estimates only need finite lists of rational input points
and the finitely many integer breakpoints.

\paragraph{From the mean back to primes.}
The floor identity and the lemma for $g_1$ give $m(x)\to0$.
Define $L_x=\sum_{d\le x}\mu(d)\log(x/d)/d$. Grouping by $dk$ gives
\[
 1=\sum_{d\le x}\frac{\mu(d)}d H_{\lfloor x/d\rfloor}
  =L_x+\gamma m(x)+x^{-1}\sum_{d\le x}\mu(d)g_2(x/d).
\]
Hence $L_x\to1$. Finally $\log n=\sum_{r\mid n}\Lambda(r)$, and
inversion, then summation over $n\le x$, gives
$\psi(x)=\sum_{d\le x}\mu(d)\log(\lfloor x/d\rfloor!)$.
Using $g_3$ gives
$\psi(x)/x=L_x-m(x)+x^{-1}\sum_{d\le x}\mu(d)g_3(x/d)\to1$.
A modulus for $m(x)\to0$ alone gives $M(x)/x\to0$ by
\[
 M(x)=xm(x)-\sum_{j=1}^{x-1}m(j).
\]
Indeed if $|m(j)|\le\eta$ for $j\ge U$, dividing by $x$ costs at
most $2\eta+2U/x$, using the global bound $|m(j)|\le2$.

\paragraph{The prime-counting formulation.}
Put $\vartheta(x)=\sum_{p\le x}\log p$. Prime powers of exponent
at least two contribute at most $\sqrt x\log x$ to $\psi(x)$:
each prime $p\le\sqrt x$ contributes at most $\log x$.
For $r\ge2$ put $y=x^{1-1/r}$ (the cutoff can be compared by integer
$r$th powers). Splitting primes at $y$ gives the elementary bounds
\[
 (1-1/r)\big(\pi(x)\log x-y\log x\big)
       \le\vartheta(x)\le \pi(x)\log x .
\]
The error $y\log x/x$ is at most $2r x^{-1/(2r)}$ since
$\log u\le u$ for $u\ge1$. Choosing $1/r$ below a quarter of the desired
error and $x$ above a suitable integer power of $2r/\epsilon$
gives each implication between the normalized $\vartheta$ and
$\pi$ limits. Taking a finer fixed fraction of the target precision
absorbs the factor $(1-1/r)^{-1}\le2$. These are PR modulus
transformations. The same power bound controls $\log x/\sqrt x$.

\section{An elementary transfer from prime counts to signed cancellation}
\label{sec:bootstrap}
There is a second route to the boundary estimates. It uses only finite
summation identities after importing a quantitative prime number theorem.
Write $L(t)=1+\lfloor\log_2 t\rfloor$ for the binary length of an integer
$t\ge1$, and let $\Lambda(p^j)=\log p$ for a prime $p$ and $j\ge1$,
with value zero elsewhere. Fix any integer $A\ge1$. The required input is
\begin{equation}\label{eq:bootstrap-input}
    |E(t)|\le Kt/L(t)^A,\qquad E(t)=\sum_{m\le t}(\Lambda(m)-1).
\end{equation}
Logs here denote rational enclosure names; the multiplier $K$ and every
binary length are integers.

\paragraph{Obtaining the input from ProveIt.}
The theorem \texttt{thm:theta-modulus-established} in the repository's
\texttt{pnt\_li\_ipra\_2026-11-15/tex/theta\_bridge.tex} gives, with
$C=2^{240}$ and $\vartheta(t)=\sum_{p\le t}\log p$,
\[
  t\ge2^{C(r+1)^{10}}\quad\Longrightarrow\quad
  |\vartheta(t)-t|\le2^{-r}t.
\]
It is the upstream quantitative theorem used in this section.
For a fully specified instance of \eqref{eq:bootstrap-input} take
\[
 D_A=2(4A+4)^{A+1},\qquad
 K=(C+2)^{A+1}+4C^A(20A)^{10A}+D_A.
\]
Indeed set $L=L(t)$ and choose the largest $k$ with $Ck^{10}\le L-1$
by bounded search. For $k=0$ use
$|\vartheta(t)-t|/t\le1+\log t\le1+L$ and $L\le C$.
For $k\ge1$ the theorem bounds this ratio by $2^{1-k}$, while
$L\le C(k+1)^{10}$. The power-series inequality
$e^u\ge u^d/d!$ and $\log2\ge1/2$ give
$(k+1)^d 2^{-k}\le 2(2d)^d$ for $d\ge1$.
This bounds the $\vartheta$ contribution by the first two terms of $K$.
For each prime $p\le\sqrt t$, its powers of exponent at least two contribute
in total at most $\log t$. Thus
$0\le \sum_{m\le t}\Lambda(m)-\vartheta(t)\le\sqrt t\log t$.
Since $t\ge2^{L-1}$, the same exponential inequality gives
$L^{A+1}/\sqrt t\le D_A$, as required.

\paragraph{The finite identity.}
Let $e(m)=\Lambda(m)-1$. The operator $Df(n)=f(n)\log n$ is a derivation
for Dirichlet convolution, by $\log(ab)=\log a+\log b$. Write $\mathbf1$
for the constant-one function. Factorization gives
$D\mathbf1=\mathbf1*\Lambda$. Differentiate $\mu*\mathbf1=[n=1]$
and convolve on the right with $\mu$ to obtain $D\mu=-\mu*\Lambda$.
Summing this equality for $n\le x$ and using
$\sum_{dm\le x}\mu(d)=1$ gives for each positive integer $x$
\begin{equation}\label{eq:bootstrap-identity}
 M(x)\log x
 =\sum_{d\le x}\mu(d)\log(x/d)-1
       -\sum_{dm\le x}\mu(d)e(m).
\end{equation}
All sums on the right are finite.

\paragraph{Gain of a half-power in one step.}
Set $B_0=1$ and recursively take
\[
 B_{j+1}=16^{A+1}+64\cdot4^A(K+2)(B_j+D_A+1).
\]
We prove $|M(x)|\le B_j x/L(x)^{j/2}$ for $0\le j\le A$. The case
$j=0$ uses $|\mu|\le1$. Put $s=j/2$, $L=L(x)$, and suppose $j<A$.
For $L<16$ use the trivial bound and $B_{j+1}\ge16^{A+1}$. Otherwise
choose $U=2^{\lfloor\sqrt L\rfloor}$. In the double sum of
\eqref{eq:bootstrap-identity}, the part $m\le U$ is
$\sum_{m\le U}e(m)M(\lfloor x/m\rfloor)$. The binary length of
$\lfloor x/m\rfloor$ is at least $L-\lfloor\sqrt L\rfloor\ge L/2$.
A dyadic block $2^h\le m<2^{h+1}$ contributes at most $2(K+2)$
to $\sum |e(m)|/m$, since \eqref{eq:bootstrap-input} gives
$\sum_{m\le y}|e(m)|\le(K+2)y$. Thus the small-$m$ part has magnitude
at most
\[
       4(K+2) 2^s B_j x\, L^{1/2-s}.
\]
The remaining part is exactly
$\sum_{d\le x/U}\mu(d)(E(\lfloor x/d\rfloor)-E(U))$.
Each argument of $E$ is at least $U$, with binary length at least
$\sqrt L$. Since $\sum_{d\le x}1/d\le 1+\log x\le1+L$, its magnitude
is at most $3Kx L^{1-A/2}$.

Finally summation by parts rewrites the first sum on the right of
\eqref{eq:bootstrap-identity} as
$\sum_{d<x} M(d)\log(1+1/d)$. For $d<\sqrt x$ bound
$|M(d)|\le d$; for $d\ge\sqrt x$ use the inductive bound and
$L(d)\ge L/4$. The result is bounded by
$\sqrt x+4^s B_j x/L^s$.
Use $L^s/\sqrt x\le D_A$ and divide \eqref{eq:bootstrap-identity} by
$\log x\ge L/4$. As $s+1/2\le A/2$, the three estimates together
are at most $B_{j+1}x/L^{s+1/2}$, proving the step. All block endpoints
come from integer roots or bit lengths. The fractional powers in the
bounds for $M$ disappear on squaring nonnegative quantities. For
quantifier-free induction one can fix an arbitrary upper bound $X$ and
carry the Boolean conjunction of the integer estimates for $1\le x\le X$
at stage $j$; every use of the inductive hypothesis is within this list.

Taking $A=2b$ proves an effective $|M(x)|\le B_{2b}x/L(x)^b$ for
each positive integer $b$. Thus the repository's explicit prime estimate
supplies whatever fixed logarithmic saving is needed in ordinary
logarithmically weighted sums of $\mu(n)/n$.

\section{Lean formalization and object-theory endpoints}
\label{sec:lean-route}

\subsection{Two independently usable representations}
The file \texttt{lean/FiniteInversion.lean} uses only Lean's bundled
library. It was checked with Lean 4.32.0, the version pinned by ProveIt,
and with Lean 4.19.0. Entries of its matrix are genuine Lean integers.
If $W(i,j)$ is the weight at row $i$, column $j$, write
\[
 (Lf)(i)=\sum_{j<i}W(i,j)f(j),\qquad (Zf)(i)=f(i)+(Lf)(i).
\]
Define $b_0(i)=0$, $b_{t+1}(i)=g(i)-(Lb_t)(i)$, and $b(i)=b_{i+1}(i)$.
If $i<t$, induction on $t$ gives $b_{t+1}(i)=b_t(i)$, since every
previous entry needed in row $i$ has index $j<i$. Thus $Zb=g$, and
substitution, from the lowest row upwards, gives $b=f$ whenever $g=Zf$.
These are \texttt{integrate\_invert} and \texttt{invert\_integrate}.
The file also proves uniqueness on a specified prefix of rows.

Subtract two consecutive iterates. The first difference is $g$, and
each remaining difference is $-L$ applied to its predecessor. Hence
\[
 b(i)=\sum_{r=0}^{i} ((-L)^r g)(i).
\]
Each summand is a signed sum over strictly increasing length-$r$ chains,
by expanding its last edge. For $r>i$ the row vanishes.
The file checks both assertions uniformly in $i,r,W,g$.
With $W(i,j)=[j+1\mid i+1]$ and $g(i)=[i=0]$, the recurrence computes
$\mu(i+1)$. Its identification with the prime-factor definition follows
from the divisor cancellation proof and uniqueness. The Lean file proves
divisor cancellation and the chain formula for its recurrence-defined
function. A second proof, \texttt{ExponentBox.box\_cancellation}, sums the
prime-factor definition over a list of exponent ranges of arbitrary
length. The coordinate factor at $e=0,1,\ge2$ is respectively $1,-1,0$.
Its prefix sum up to exponent $a$ is $[a=0]$. Factoring this sum off
the first coordinate and iterating proves the product formula. The human
divisor--exponent bijection connects this calculation with the numerical
divisor scan; the standalone file leaves that bijection to the report.

The companion \texttt{FormalSeries.lean} imports this file. The coefficient
calculations in Section~\ref{sec:mob-basel} describe its scaled
representations and conclusions. Run \texttt{sh lean/check.sh} to check
both files together.

\subsection{The existing internal proofs}
For object proofs use revision \texttt{65b76a5b8aa8} of ProveIt and the
namespace \texttt{LeanProofs.}\allowbreak\texttt{IPRA.}\allowbreak\texttt{AnalyticNumberTheory}. The relevant
source directory is
\href{https://github.com/VladimirReshetnikov/ProveIt/tree/65b76a5b8aa87ecad4103c2d287eff5660d3c84c/Logic/PrimitiveRecursiveArithmetic/Derivations/Lean/IPRADerivations/AnalyticNumberTheory}{the analytic number theory derivations directory}.
The ordinary calculus has function codes \texttt{PR k}, terms
\texttt{PTm n} over $n$ variables, formulas \texttt{Fm PTm n}, and
deductions \texttt{Prf PRAAx} indexed by their contexts and conclusions.
A term body with free variables for its input and parameters expresses a
schematic arithmetic function. Object theorems use numerical variables
for bounds. They may be specialized at arbitrary terms in any context
which supplies the stated hypotheses.

The repository takes $\mu_N^+(x)=[s_N(x)=0][\omega_N(x)\text{ even}]$ and
$\mu_N^-(x)=[s_N(x)=0][\omega_N(x)\text{ odd}]$, scanning primes $q<N$.
Here $\omega_N$ counts nonzero valuations, and $s_N$ counts those at
least two. Once $N>x\ge1$ the pair is stable. Integer equations are
written by transferring negative components across the equality.
The following are source-inspected object-theory declarations at that
revision. The replay file lists their fully resolved module imports.
\begin{itemize}
\item \texttt{Mobius.swap}, \texttt{Mobius.kill},
\texttt{Mobius.mobius\_one}, and \texttt{Mobius.mobius\_cancel} give
the prime steps and the divisor sum. \texttt{MobiusBounds.cancel\_gen}
allows an arbitrary prime bound and an arbitrary divisor scan limit,
both strictly above the positive argument. This matters when one cutoff
is shared by many terms.
\item \texttt{MobiusInversion.mobius\_inversion} instantiates the divisor
inversion at a term body and \texttt{phi\_mobius} derives its totient
specialization. For signed input bodies use the theorem on each natural
component and add the cross terms.
\item \texttt{SummatoryPhi.summatory\_phi} gives the floor-square
totient identity. \texttt{MuLambda.mu\_lambda} gives
$\mu*\Lambda_w=-\mu L_w$ for the prime weight represented by the
given body. \texttt{PsiModulus.}\allowbreak\texttt{psi\_modulus} supplies the quantitative
PNT for the repository's finite harmonic weights; \texttt{FinalPNT.}\allowbreak\texttt{arith\_pnt}
gives its prime-counting version.
\end{itemize}
These existing derivations and the new standalone Lean proofs address
different judgments. We inspected the pinned upstream sources; this
packet includes replay commands for a prepared checkout. We ran both
standalone files locally. The remaining object derivations described below
are a formalization plan.

\subsection{Constructing additional internal deductions}
\begin{enumerate}
\item \emph{Encoding and induction.} Use \texttt{PR.comp} to compose
body codes and \texttt{PR.prec} with arguments in the order
(counter, parameters), its step receiving (predecessor, previous value,
parameters). Use bounded sums and products from \texttt{Core/Definitions}
and \texttt{Core/MultiSums}, \texttt{Core/Search} for the first matching
index, and \texttt{Core/Fubini} for two-index rearrangements. Define
equality of signed pairs by $u_++v_-=v_++u_-$. For a bounded family of
integer equalities $E_i$, form its failure count
$\sum_{i<t}(1-[E_i])$. Induct on its equality to zero using
\texttt{PRA.qfInd}. Nested counts allow several bounded indices.
Recover each equality since a nonzero summand contributes at least one.
This implements every use of finite strong induction on the prefix.
\item \emph{Inversion and posets.} Fix a packed finite relation and a
numerical bound $B$. On $i<B$ implement $L$ by the displayed bounded
sum. Store the $B$ entries of the current row vector as a finite list
of signed pairs; compute the next vector by one length-$B$ fold.
The computation of $b_{i+1}(i)$ now recurses on a natural code at fixed
$B$ instead of iterating over functions. Induct on the zero failure count
of the equalities $b_{t+1}(j)=b_t(j)$ for $j<\min(B,t)$.
Base, step and uniqueness are the finite proofs above. For an interval
$[a,b]$ in a poset use a finite list with $x<y$ implying earlier
position for $x$. Set $W(i,j)=[x_j<x_i]$ and $g(i)=[x_i=a]$.
On divisibility take $B=n$, $x_i=i+1$, and instantiate
\texttt{MobiusBounds.cancel\_gen} to compare the two definitions of
$\mu$ pointwise by prefix uniqueness.
\item \emph{Coefficient identities.} Store degrees $0$ through $B$,
representing a polynomial by a signed coefficient list. Convolution at
$k\le B$ is the fold $\sum_{j\le k}a_j b_{k-j}$. For Dirichlet
products replace the additive incidence by $[de=k]$ and scan only
$1\le d,e\le B$. For the Lambert series the coefficient of $X^k$
adds one term per $d$ with $d\mid k$. Apply
\texttt{MobiusBounds.cancel\_gen} at that $k$ and close all degrees
through $B$ by the zero failure count. Formal logarithms may be stored
with common denominator $(B+1)!$; prove divisibility before exact
division. Polynomial cyclotomic arguments use lists with the explicitly
given degree bound, monic division with remainder, and equality of
coefficients. Finite geometric sums are first proved by recursion on
their length, so evaluating at a unit in a specified quotient ring
uses only addition and multiplication of integer pairs or lists.
For the irreducible count over a field of cardinality $q$, supply two
$q$-by-$q$ operation tables and validate the field axioms by a bounded
scan. Represent a monic of degree $n$ by the base-$q$ digits of a number
less than $q^n$. For $n>0$ its irreducibility bit tests every monic
divisor of degree $1$ through $\lfloor n/2\rfloor$ by remainder.
Summing this bit gives $I_q(n)$. If a division splits a polynomial of
positive degree, both new degrees decrease. Recursion through a coded
list, with sum of the pending degrees bounded by $n$, gives a complete
factor list in at most $2n-1$ node visits. Use the Euclidean B\'ezout
identity to prove Euclid's lemma for monic irreducibles, then compare
multiplicities. This supplies the factorization bijection used by the
finite Euler product.
\item \emph{Finite fields and Fourier sums.} Encode the two $q$-by-$q$
field operations by lists of entries $<q$. Test associativity, distributivity
and the inverse laws on the bounded lists of tuples. Division by a monic
polynomial is a fold on its degree, stripping the leading term. Record
quotient and remainder and verify the product identity at each step.
For degree $d$ enumerate coefficient codes $<q^d$ and enumerate the
possible monic factors through degree $\lfloor d/2\rfloor$; the bit
rejecting all of them defines $I_q(d)$. Induct on degree to establish
the unique factor multiset, then count products in the coefficient table
through degree $n$ and apply the existing inversion theorem to the term
$d I_q(d)$. For $C_n(h)$, supply the Euclidean-algorithm certificate
$(U,V)$ for each invertible difference. In the quotient ring the checker
verifies its inverse by remainder equality. The same geometric-sum proof
over the table of residues modulo $p$ gives the congruence.
\item \emph{Approximation statements.} A rational name has primitive
recursive numerator and positive denominator at each requested precision.
Implement a complex name by two such pairs. For a fixed input code or a
specified PR family and precision $k$, substitute the displayed cutoff in
the relevant finite sum or product, perform exact denominator clearing,
and express the error as two integer inequalities per real coordinate.
For a tail use an additional free end index $V\ge U$, and prove the
displayed estimate for every such $V$ by induction on $V-U$. Pass
between different cutoffs by the triangle inequality and the given
moduli. At the PNT step instantiate \texttt{psi\_modulus} with the
repository's own mesh and denominator; compare alternative rational
meshes by refining to their product. These operations leave all
induction formulas quantifier-free after the bounded checks are compiled.
A statement about a family of names quantifies over indices and precisions
for that fixed family.
\item \emph{The strong summatory estimate and boundary moments.}
The ordinary PNT modulus suffices for Section~\ref{sec:pnt}.
For Theorem~\ref{thm:strongM} the stronger analytic input in \cite{strip}
must itself be represented by finite deductions. Implement the Euler
summation samples with rational logarithm and exponential enclosures,
the fourth-power strip comparison, and the regular reciprocal evaluation
at the pole. The finite contour invariant equates the signed internal
edge sum to zero, then accumulates at most one Taylor error per cell.
Apply it to the smoothed signed sum at the two rational arguments.
After raising the unsmoothing estimate to an eighth power, the closed
matrix compares $2^k|M(x)|^8$ and $x^8$. To remove arbitrary precision
from an inequality between integers it suffices to make the additive
enclosure slack less than $1$: discreteness then gives the exact weak
inequality. The moment target is the rational partial-sum comparison
following \eqref{eq:momentTail}. A second possible dependency is the
strong $\vartheta$ modulus and the bounded numerical induction in
Section~\ref{sec:bootstrap}. Both routes require additional embedded
deductions.
\end{enumerate}

\section{Coverage of the source articles}\label{sec:coverage}
The table follows MathWorld's equation numbers \cite{mw}. The numbered
claims are all covered. The entries on the right identify the arithmetic
content of the corresponding infinite formulas.

\begin{center}\small
\begin{tabular}{@{}p{.13\linewidth}p{.40\linewidth}p{.40\linewidth}@{}}
\toprule
Equations & Statement & Proof or interpretation\\\midrule
(1), (2) & Definition of $\mu$, summatory $M$ &
Valuations and bounded sums, \S\ref{sec:core}.\\
(3), (4) & Reciprocal Dirichlet series and Euler factors &
Positive half-plane margin, prime-mask expansion, \S\ref{sec:dirichlet-interior}.\\
(5), (12) & Lambert series and the exponential product &
Truncated coefficients and analytic error bounds, \S\ref{sec:mob-lambert}.\\
(6), (7) & The first two moments at $s=1$ &
Ordinary partial sums with a PR modulus, \S\ref{sec:boundary}.\\
(8)--(10) & The three reciprocal-square masses &
Finite trigonometric proof of the two even zeta values, \S\ref{sec:mob-basel}.\\
(11), (13) & Squarefree divisor count and product law &
Prime masks and additive valuations, \S\ref{sec:core}.\\
(14), (15) & Divisor cancellation and $\mu*\tau=\mathbf1$ &
The involution and convolution associativity, \S\ref{sec:core}.\\
\bottomrule
\end{tabular}
\end{center}

The additional formulas on Wikipedia \cite{wiki} include the Liouville
description, $\zeta(s)/\zeta(2s)$, the prime-shifted Lambert series,
the second logarithmic moment at $s=1$, the trigonometric unital sum
\cite{kline}, the two primitive-root identities, the Farey sum, the
irreducible-polynomial count, and convolution powers with a polynomial
parameter. Sections~\ref{sec:core}--\ref{sec:pnt} prove these and the
mean-value equivalence with PNT. Section~\ref{sec:boundary} also proves a conditional zero-exclusion bound from power cancellation of $M$. A computable finite-interval presentation
is specified for the incidence-algebra generalization. The occupation
number argument explains the arithmetic content of the gas model.

The value lists, including all sixty tabulated values, the initial sign
fibers, and the displayed products of three and five distinct primes,
are reproduced by \texttt{checks.py}. For example, bounding the scan
by 9690 certifies completeness through the last listed product of five
primes. The script also checks independent sample counts over two finite
fields. The published plots are illustrations of the same decidable
tests. The articles' historical notes, software listings, and references
to other problems have no additional theorem statement about $\mu$ to
prove here. Comparing the complexity of the two primitive-root and
prime-factor evaluation procedures would require a specified computational
model and representation of the output. We give exact bounded algorithms
for both and make no complexity comparison.

\begin{thebibliography}{9}\raggedright
\bibitem{wiki} Wikipedia contributors, \emph{Möbius function}, article
used for the formula inventory. \href{https://en.wikipedia.org/wiki/M\%C3\%B6bius_function}{Article}.
\bibitem{mw} E. W. Weisstein, \emph{Möbius Function}, MathWorld,
equations (1)--(15). \href{https://mathworld.wolfram.com/MoebiusFunction.html}{Article}.
\bibitem{kline} J. Kline, \emph{Unital Sums of the Möbius and Mertens Functions},
Journal of Integer Sequences 23 (2020), article 20.8.1.
\href{https://cs.uwaterloo.ca/journals/JIS/VOL23/Kline/kline4.html}{Paper}.
\bibitem{proveit} V. Reshetnikov et al., \emph{ProveIt},
\texttt{Logic/PrimitiveRecursiveArithmetic}, in particular
\texttt{ConciseIPRA} and \texttt{Derivations}. Declarations inspected at
\texttt{65b76a5b8aa8} (full revision in the link).
\href{https://github.com/VladimirReshetnikov/ProveIt/tree/65b76a5b8aa87ecad4103c2d287eff5660d3c84c/Logic/PrimitiveRecursiveArithmetic}{Pinned directory}.
\bibitem{strip} ProveIt, \emph{Logarithmic-integral strengthenings of the
prime number theorem in iPRA}, November 15, 2026.
The files \texttt{tex/analytic.tex} (E1)--(E8), (A1)--(A7),
\texttt{tex/finite\_contour.tex}, and \texttt{tex/theta\_bridge.tex}
supply the quantitative analytic lemmas reused in
Sections~\ref{sec:boundary} and~\ref{sec:bootstrap}.
These are human proof sources. The new signed smoothing and the boundary
moment deductions are given in this report.
\href{https://github.com/VladimirReshetnikov/ProveIt/tree/a9c0f1949be61679ed2e219b97dc42ecb7ae91a9/Logic/PrimitiveRecursiveArithmetic/docs/proofs/analytic-number-theory/pnt_li_ipra_2026-11-15}{Pinned packet}.
\end{thebibliography}

\end{document}
